\section{Uniform Convergence}

\begin{definition}[uniform convergence]
\mymark{***}
Let $A$ be a non-empty set and
$f\colon A \to \RR$.
We define the \emph{uniform norm}%
\mymargin{uniform norm}\index{norm!uniform} (or $\sup$ norm)
of $f$ as
\[
  \Abs{f}_\infty = \sup_{x\in A} \abs{f(x)}
\]

If $g\colon A \to \RR$ as well,
we define the \emph{uniform distance}
\mymargin{uniform distance}
\index{distance!uniform}
between $f$ and $g$ as
\[
  d_\infty(f,g) = \sup_{x\in A} \abs{f(x)-g(x)}.
\]

If $f_k$ is a sequence of functions and $f$ is a function, we say that $f_k$
\emph{converges uniformly}
\mymargin{uniform convergence}%
\index{convergence!uniform}%
to $f$
and we write
\[
f_k \To f
\] if
$d_\infty(f_k,f)\to 0$.
\end{definition}


\begin{example}
\label{ex:466533}
The sequence
\[
f_k(x) = \sqrt{x^2 + \frac{1}{k}}
\]
converges uniformly (on all of $\RR$) to the function $f(x) = \abs{x}$. Indeed,
setting $g_k(x) = f_k(x) - f(x)$,
the function $g_k$ is differentiable for $x\neq 0$ and for $x>0$ we have
\[
  g'_k(x) = \frac{x}{\sqrt{x^2+\frac 1 k}} - 1 < 0.
\]
Thus, the function $g_k$ is decreasing on $[0,+\infty)$. By symmetry (it is an
even function), it is increasing on $(-\infty, 0]$.
Therefore, the maximum of $g_k$ is at $x=0$. Clearly, $g_k \ge 0$, so we have:
\[
  \Abs{f_k - f}_\infty = \sup_{x\in \RR} g_k(x) = g_k(0) = \frac{1}{\sqrt k} \to 0.
\]
Thus, $f_k \To f$.
\end{example}

Observe that in general $\Abs{f}_\infty$ and $d_\infty(f,g)$ can take the value
$+\infty$ (for example, if $A=\RR$, $f(x)=x$ and $g(x)=0$)
and therefore they are not necessarily
a norm and a distance.

\begin{theorem}[properties of the uniform norm]
The uniform norm satisfies all the properties of a norm
(Definition~\ref{def:norma}), except that it can take values in $[0,+\infty]$
instead of $[0,+\infty)$.
\end{theorem}
%
\begin{proof}
Clearly, the uniform norm does not take negative values as it is the supremum of
a (non-empty) set of non-negative real numbers. Moreover, if $\Abs{f}_\infty=0$,
it means that $\abs{f(x)}=0$ for every $x$, and thus $f=0$ (separation
property).

Homogeneity follows from the homogeneity of the absolute value, as we have
\[
  \sup_{x\in A} \abs{(\lambda \cdot f)(x)}
  = \sup_{x\in A}\abs{\lambda \cdot f(x)}
  = \sup_{x\in A}\abs{\lambda}\cdot \abs{f(x)}
  = \abs{\lambda} \cdot \sup_{x\in A}\abs{f(x)}.
\]

The triangle inequality follows from the triangle inequality of the absolute
value, which is preserved when taking the supremum:
\[
  \sup_{x\in A} \abs{f(x)+g(x)}
  \le \sup_{x\in A} \Enclose{\abs{f(x)} + \abs{g(x)}}
  \le \sup_{x\in A} \abs{f(x)} + \sup_{x\in A} \abs{g(x)}.
\]
\end{proof}

\begin{theorem}[completeness of bounded functions]%
  \label{th:limitate_completo}%
Let $A$ be a set.
The vector space
of bounded functions $f\colon A \to \RR$
(that is, functions with finite uniform norm)
\[
  \B(A) = \ENCLOSE{f\in \RR^A\colon \Abs{f}_\infty < +\infty }
\]
endowed with the uniform norm $\Abs{\cdot}_\infty$ is a Banach space (i.e., a
normed and complete vector space).
On this Banach space, the distance induced by the norm is the uniform distance
$d_\infty$, and the convergence induced by the distance is uniform convergence.
\end{theorem}
%
\begin{proof}
By definition, it is verified that the uniform norm $\Abs{\cdot}_\infty$ takes
finite values on $\B(A)$.
Thus, according to the previous theorem, $\Abs{\cdot}_\infty$ is indeed a norm,
and $\B(A)$ is therefore a normed space. We now demonstrate that it is complete,
i.e., that Cauchy sequences converge.

Let $f_k$ be a Cauchy sequence in $\B(A)$.
Then for every $x\in A$, $f_k(x)$ is a Cauchy sequence in $\RR$ since (by
definition of $\sup$)
\[
  \abs{f_k(x) - f_j(x)} \le \Abs{f_k - f_j}_\infty
\]
and thus if $\Abs{f_k- f_j}_\infty < \eps$
it follows that for fixed $x\in A$, $\abs{f_k(x)-f_j(x)} < \eps$.

Therefore, for every $x\in A$, the numerical sequence $f_k(x)$ converges since
$\RR$ is complete. Setting $f(x) = \lim f_k(x)$, we have thus found a candidate
limit for the sequence.
We must now show that $f\in \B(A)$ and that $f_k$ converges uniformly to $f$.
For every $\eps>0$, by the Cauchy condition, there must exist $N\in \NN$ such
that if $k,j>N$, then
\[
  d_\infty(f_k,f_j) < \eps.
\]
But then for every $x\in A$, for every $k>N$, and for every $j>N$, we have:
\[
  \abs{f_k(x) - f(x)} \le \abs{f_k(x) - f_j(x)} +
  \abs{f_j(x) - f(x)} < \eps + \abs{f_j(x)-f(x)}.
\]
Since for every $x$, $f_j(x) \to f(x)$, for every $x$, there exists a $j$ such
that $\abs{f_j(x)-f(x)} < \eps$, and thus we can conclude that
\[
  \abs{f_k(x)-f(x)} < 2\eps.
\]
Taking the $\sup$ for $x\in A$, we obtain
\[
  \Abs{f_k -f}_\infty \le 2 \eps.
\]
We have thus verified the definition of limit $\Abs{f_k -f}_\infty\to 0$. In
particular, $\Abs{f}_\infty < +\infty$ since the triangle inequality holds
\[
  \Abs{f}_\infty \le \Abs{f-f_k}_\infty + \Abs{f_k}_\infty < +\infty
\]
as $\Abs{f-f_k}_\infty \to 0$ and $\Abs{f_k}_\infty < +\infty$.
\end{proof}

\begin{definition}[pointwise convergence]
\mymark{***}
Let $f_k\colon A \to \RR$ be a sequence of functions
and let $f\colon A \to \RR$ be a function.
If for every $x\in A$, $f_k(x)\to f(x)$, we say that
the sequence $f_k$
\emph{converges pointwise}
\mymargin{pointwise convergence}
\index{convergence!pointwise}
to $f$.
\end{definition}

\begin{theorem}[uniform convergence implies pointwise convergence]
\mymark{***}
Let $f_k\colon A \to \RR$ be a sequence of functions.
If $f_k$ converges uniformly to a function $f$, then $f_k$ converges pointwise
to $f$.
\end{theorem}
%
\begin{proof}
It is sufficient to observe that for every $x\in A$, we have
\[
  \abs{f_k(x)-f(x)} \le \sup_{y\in A} \abs{f_k(y)-f(y)}
   = \Abs{f_k-f}_\infty \to 0.
\]
\end{proof}

\begin{example}[sequence that converges pointwise but not uniformly]
\mymark{***}%
\label{ex:puntuale_non_uniforme}%
Let $f_k\colon [0,1]\to \RR$ be the sequence of functions defined by
$f_k(x)=x^k$. If $x\in[0,1)$, we have $x^k \to 0$, while if $x=1$, we have $x^k
\to 1$.
Thus, the sequence $f_k$ converges pointwise to the function
\[
f(x) =
 \begin{cases}
  0 & \text{if $x\in [0,1)$}\\
  1 & \text{if $x=1$}.
 \end{cases}
\]
However, observe that
\[
  d_\infty(f_k,f) = \sup_{x\in [0,1]} \abs{f_k(x)-f(x)}
  \ge \lim_{x\to 1^-} \abs{f_k(x) - f(x)} = 1.
\]
thus there cannot be uniform convergence of $f_k$ to $f$.
\end{example}

It is easy to convince oneself that the sequence $f_k$ in the previous example,
besides not converging uniformly, does not admit any uniformly convergent
subsequence. Therefore, such a sequence cannot be contained in any compact
subset of $C^0([0,1])$. In particular, the unit disk
\[
  D = \ENCLOSE{f\in C^0([0,1])\colon \Abs{f}_\infty \le 1}
\]
is a closed and bounded set that is not compact.

\begin{theorem}[continuity of the uniform limit]
\mymark{***}%
Let $X$ be a metric space and let $f_k\colon X\to \RR$
be continuous functions that
converge uniformly to a function $f\colon X \to \RR$. 
Then $f$ is also continuous.
\end{theorem}
%
\begin{proof}
\mymark{***}
Fixing $x_0\in X$, it suffices to show that for every $\eps>0$
there exists $\delta>0$ such that if $d(x,x_0)< \delta$, then $\abs{f(x)-f(x_0)}
< 3 \eps$.
By definition of uniform convergence, given $\eps>0$,
there exists an $N\in \NN$ (in fact, infinitely many) such that
$d_\infty(f_N,f)< \eps$. By the continuity of $f_N$, for the same $\eps$, there
exists $\delta>0$
such that if $d(x,x_0) < \delta$, then $\abs{f_N(x)-f_N(x_0)} < \eps$. But then
if $d(x,x_0)<\delta$, we have
\begin{align*}
\abs{f(x)-f(x_0)}
&\le \abs{f(x) - f_N(x)}
 + \abs{f_N(x)-f_N(x_0)}
 + \abs{f_N(x_0) - f(x_0)} \\
 &\le \Abs{f-f_N}_\infty + \eps + \Abs{f-f_N}_\infty
  \le 3\eps.
\end{align*}
\end{proof}

\begin{theorem}[completeness of $C^0({[a,b]})$]%
\label{th:C0_completo}
\mymark{***}%
\mymargin{$C^0([a,b])$ is complete}%
\index{completeness!of $C^0([a,b])$}%
The space $C^0([a,b])$ of continuous functions defined on a closed and bounded
interval,
endowed with the uniform norm $\Abs{\cdot}_\infty$, is a Banach space (i.e., a
normed and complete vector space).
\end{theorem}
%
\begin{proof}
By the Weierstrass theorem, every continuous function defined on the compact
$[a,b]$ is bounded.
Thus, $C^0([a,b])$ is a vector subspace of $\B([a,b])$. 
Moreover, the previous theorem (continuity of the limit) tells us that
$C^0([a,b])$ is a closed subspace of $\B([a,b])$.
But $\B([a,b])$ is complete, and therefore $C^0([a,b])$, being closed in
$\B([a,b])$, is also complete.
\end{proof}

The uniform norm is the natural norm on $C^0([a,b])$ as it makes it a complete
space. For this reason, the uniform norm on continuous functions
is also called the \emph{$C^0$ norm} and can be denoted as follows:
\index{$\Abs{\cdot}_{C^0}$}
\index{norm!$C^0$}
\[
  \Abs{f}_{C^0} = \Abs{f}_{C^0([a,b])} = \Abs{f}_\infty
  \qquad\text{for $f\in C^0([a,b])$.}
\]